
\titledquestion{Multiple Choices}

Each question has \textbf{one or more} correct answer(s).
Select all the correct answer(s). Each question is worth 3 points.
You will receive 3 points if you select all correct answers,
1 point if you select a non-empty proper subset of the correct answers,
and 0 points if you select any incorrect answer or leave the answer empty.


Write your answers in the following table.

\begin{table}[htbp]
    \centering
    \begin{tabular}{|p{1.5cm}|p{1.5cm}|p{1.5cm}|p{1.5cm}|p{1.5cm}|}
        \hline
        (a) & (b) & (c) & (d) & (e) \\
        \hline
        %%%%%%%%%%%%%%%%%%%%%%%%%%%%%%%%%%%%%%%%%%%%%%%%%%%%%%%%%%
        % YOUR ANSWER HERE.
          ADE  &  ABE   &  ACE   &   AB  &  ACE   \\
        %%%%%%%%%%%%%%%%%%%%%%%%%%%%%%%%%%%%%%%%%%%%%%%%%%%%%%%%%%
        \hline
    \end{tabular}
\end{table}

\begin{parts}

    \part[3] Consider a rooted tree with $n \geq 2$ nodes.
    Which of the following statements is(are) necessarily true?

    \begin{choices}
        \choice A node at depth $d$ has exactly $d+1$ ancestors.

        \choice Deleting one leaf at the maximum depth always
        decreases the height of the tree.

        \choice Detaching a subtree rooted at $u$ and reattaching
        $u$ as a child of one of its own strict descendants
        always produces a valid tree.

        \choice The sum of the degrees of all nodes is $n-1$.

        \choice Detaching one child subtree decreases
        its former parent's degree by exactly one.
    \end{choices}

    \part[3] Which of the following statements about binary
    trees is(are) correct?

    \begin{choices}
        \choice A complete binary tree with exactly
        $2^k$ nodes, where $k \geq 1$, has exactly one node
        at depth $k$.

        \choice There exists a full binary tree with
        exactly 23 nodes and height 11.

        \choice Every perfect binary tree is complete,
        but it need not be full.

        \choice Every full binary tree of height $h \geq 2$
        has at least $2^h$ leaves.

        \choice A complete binary tree with an even
        positive number of nodes has exactly one node with
        only one nonempty child.
    \end{choices}

    \part[3] Consider the following recursive function.
    Assume that accessing a child pointer and computing
    \texttt{max} each take constant time.

    \begin{algorithm}[H]
        \caption*{F(u)}
        \begin{algorithmic}[1]
            \Function{F}{u}
                \If{$u = \texttt{nullptr}$}
                    \State \Return $-1$
                \EndIf
                \State $a \gets F(u.\mathrm{left})$
                \State $b \gets F(u.\mathrm{right})$
                \State \Return $1+\max(a,b)$
            \EndFunction
        \end{algorithmic}
    \end{algorithm}

    Let $n \geq 1$ be the number of nodes and $h$ the height
    of the input binary tree. Which of the following
    statements is(are) correct?

    \begin{choices}
        \choice \texttt{F(root)} returns the height
        of the tree.

        \choice For every binary tree, the running time
        of \texttt{F(root)} is $O(h)$.

        \choice The maximum number of simultaneously
        active calls is $\Theta(h+1)$.

        \choice If the input is a perfect binary tree,
        the running time is $O(\log(n))$.

        \choice Replacing \texttt{return -1} with
        \texttt{return 0} increases the returned value
        by exactly one for every nonempty input tree.
    \end{choices}

    \part[3] A binary tree has the following preorder
    and postorder traversals. All node labels are distinct,
    and null pointers are not included in the sequences.

    \begin{center}
        \begin{tabular}{ll}
            Preorder:  & M, F, H, C, K, T, P, W, Y \\
            Postorder: & C, K, H, F, P, Y, W, T, M
        \end{tabular}
    \end{center}

    Which of the following could be the inorder traversal
    of this tree? Do not assume that the tree is full,
    complete, or a binary search tree.

    \begin{choices}
        \choice F, C, H, K, M, P, T, W, Y

        \choice C, H, K, F, M, P, T, Y, W

        \choice F, H, C, K, M, P, T, W, Y

        \choice F, C, H, K, M, T, P, W, Y
    \end{choices}

    \part[3] Let $T$ be a single rooted ordered tree,
    and let $B$ be its left-child right-sibling (LCRS)
    binary representation. Node labels are distinct.

    Which of the following statements is(are) correct?

    \begin{choices}
        \choice If a node's right pointer points
        to another node in $B$, those two nodes are
        siblings in $T$.

        \choice Every leaf of $T$ has both its left
        and right pointers equal to \texttt{nullptr} in $B$.

        \choice The preorder traversal of $T$
        is identical to the preorder traversal of $B$.

        \choice $T$ and $B$ always have the same height.

        \choice The postorder traversal of $T$
        is identical to the inorder traversal of $B$.
    \end{choices}

\end{parts}
